\documentclass[10pt,a4paper]{amsart}
\usepackage[margin=19mm]{geometry}
\usepackage{amsmath,amssymb,mathtools}
\usepackage[scr=rsfs,cal=euler]{mathalfa}
\usepackage{xfrac}
\usepackage[unicode,hidelinks,bookmarks=false]{hyperref}
\newcommand{\ie}{i.e.\ }
\newcommand{\cf}{cf.\ }
\newcommand{\resp}{respectively\ }
\newcommand{\Subsection}{\subsection}
\providecommand{\nobreakdash}{\nobreak\hbox{-}}
\setlength{\emergencystretch}{2em}
\multlinegap0pt
\usepackage{amssymb}
\usepackage{mathtools}
\usepackage[shortlabels]{enumitem}
\newtheorem{lemma}{Lemma}
\newcommand{\cE}{\mathcal{E}}
\title{On the spectral radius of operator tuples}
\author{Marcel Scherer, Orr Shalit,, Eli Shamovich}
\date{Source: arXiv:2605.09354v1 (CC BY 4.0)}
\begin{document}
\maketitle

\noindent\emph{Source and licence.} On the spectral radius of operator tuples. Version arXiv:2605.09354v1 (CC BY 4.0), distributed by arXiv under the Creative Commons Attribution 4.0 International licence, http://creativecommons.org/licenses/by/4.0/, as recorded in the arXiv licence metadata for this exact version and shown on the abstract page https://arxiv.org/abs/2605.09354v1. Full licence notice: \texttt{licenses/arxiv-2605.09354-CC-BY-4.0.txt}.\par\smallskip

\noindent\textit{Editorial scope.} Counted unit: the article's lemma isomorphism (source0.tex lines 923--940). The article's complete original proof of it is delivered with it (source0.tex lines 942--1018, one \texttt{proof} environment of 77 lines). No mathematics is written, repaired or re-derived: the statement and the proof are the article's own lines, in the article's own notation, and no retained line is substituted or re-flowed.  Retained uncounted as the source-local context the proof needs: lines 858--863.  Nothing was omitted from any retained span, so the package carries no omission record.  No retained line contains a citation, so the wrapper carries no bibliography and no bibliography entry.  No retained line contains a figure environment, an image-inclusion command or drawing code, so no image is delivered and the package declares no asset dependency.  Editorial formatting only, in addition: an AMS \texttt{amsart} preamble that restates the article's own declarations and macros used by the retained lines, namely the environments \texttt{lemma} on separate counters, together with the standard packages those lines need.  The retained environments print in this document as Lemma 1 in order of appearance, whereas the article prints the same items with the numbering of its own full document.  The complete native source of the article as distributed in the arXiv e-print for this exact version is bundled unchanged as \texttt{sources/200218/source0.tex}.
\begin{lemma}\label{lem: radius=norm}
    Let $\mathcal{E}$ be a selfadjoint operator space. Then, for every $X \in M_n(\cE)$ selfadjoint $\rho_{\cE}(X) = \|X\|_{\cE}$. In particular, for every $X \in M_n(\cE)$,
      \[
      \|X\|_{\mathcal{E}}=\rho_\mathcal{E}\left(\begin{pmatrix} 0 & X\\ \tau(X) & 0\end{pmatrix}\right).
      \]
\end{lemma}

\begin{lemma}\label{lem: radii linear isomorphism}
    Let $\cE_1$ and $\cE_2$ be operator spaces which are linearly isomorphic as Banach spaces, and let
  \[
    u:\cE_1\to \cE_2
  \]
be a linear isomorphism. Suppose that $\cE_2$ is selfadjoint and that there exists
$c\in M_n(\cE_2)$ such that
  \[
    \|c\|_{\max(\cE_2)}
    >
    \|u\|\,\|u^{-1}\|\,
    \|c\|_{\min(\cE_2)}.
  \]
Then there exists $X\in M_{2n}(\cE_1)$ such that
  \[
    \rho_{\min(\cE_1)}(X)<\rho_{\max(\cE_1)}(X).
  \]
\end{lemma}

\begin{proof}
Set
  \[
    \kappa=\|u\|\,\|u^{-1}\|.
  \]
Since $\cE_2$ is selfadjoint, we can define
  \[
    a=
    \begin{pmatrix}
    0 & c\\
    c^* & 0
    \end{pmatrix}
    \in M_{2n}(\cE_2).
  \]
By Lemma \ref{lem: radius=norm},
  \[
    \rho_{\min(\cE_2)}(a)
    =
    \|a\|_{\min(\cE_2)}
    =
    \|c\|_{\min(\cE_2)},
  \]
and
  \[
    \rho_{\max(\cE_2)}(a)
    =
    \|c\|_{\max(\cE_2)}.
  \]
Hence, by assumption,
  \[
    \rho_{\max(\cE_2)}(a)
    >
    \kappa\,\rho_{\min(\cE_2)}(a).
  \]
Now put
  \[
    X=u^{-1}(a)\in M_{2n}(\cE_1).
  \]
We know that the spectral radius decreases under completely contractive maps, and for a completely bounded map $T:\cE_1\to \cE_2$ it follows that
  \[
    \rho_{\cE_2}(T(x))
    \leq
    \|T\|_{\mathrm{cb}}\rho_{\cE_1}(x).
  \]
For minimal and maximal operator space structures one has
  \[
    \|u:\min(\cE_1)\to \min(\cE_2)\|_{\mathrm{cb}}=\|u\|,
    \qquad
    \|u:\max(\cE_1)\to \max(\cE_2)\|_{\mathrm{cb}}=\|u\|,
  \]
and similarly for $u^{-1}$. 
Therefore, 
  \[
    \rho_{\min(\cE_1)}(X)
    \leq
    \|u^{-1}\|\,\rho_{\min(\cE_2)}(a),
  \]
whereas
  \[
    \rho_{\max(\cE_1)}(X)
    \geq
    \frac{1}{\|u\|}\rho_{\max(\cE_2)}(a).
  \]
Since
  \[
    \rho_{\max(\cE_2)}(a)   
    >
    \|u\|\,\|u^{-1}\|\rho_{\min(\cE_2)}(a),
  \]
we obtain
  \[
    \rho_{\max(\cE_1)}(X)
    >
    \rho_{\min(\cE_1)}(X).
  \]
This completes the proof.
\end{proof}

\end{document}
